\documentclass{article}
\usepackage[UTF8]{ctex}
\usepackage{amsmath,amssymb,amsfonts}
\usepackage{enumitem}
\usepackage{geometry}
\geometry{a4paper, margin=1in}
% === 错题本排版优化（不参与题目标识） ===
\setlength{\parindent}{0pt}
\setlength{\parskip}{0.35em plus 0.10em minus 0.08em}
\setlength{\emergencystretch}{3em}
\setlist{itemsep=0.25em, topsep=0.40em, parsep=0pt}
\allowdisplaybreaks[2]
% === 错题本排版优化结束 ===
\begin{document}
\section*{国防科技大学2024年数学分析真题}
\section*{150．国防科技大学2024年数学分析与高等代数试卷}

\subsection{题目1}
1．给定两正数 \(a_{1}, b_{1}, a_{1}>b_{1}\) 。作出其等差中项 \(a_{2}=\frac{a_{1}+b_{1}}{2}\) 与等比中项 \(b_{2}=\sqrt{a_{1} b_{1}}\) 。一般地，定义

\[
a_{n+1}=\frac{a_{n}+b_{n}}{2}, b_{n+1}=\sqrt{a_{n} b_{n}}, n=1,2, \cdots .
\]

证明： \(\lim _{n \rightarrow \infty} a_{n}\) 与 \(\lim _{n \rightarrow \infty} b_{n}\) 存在且相等．

\subsection{题目2}
2．设 \(\sum_{n=1}^{\infty} a_{n}\) 为发散的非负项级数，\(S_{n}=\sum_{k=1}^{n} a_{k}\) 且 \(a_{1}>0\) 。证明：级数 \(\sum_{n=1}^{\infty} \frac{a_{n}}{S_{n}^{2}}\) 收敛。

\subsection{题目3}
3．证明：当 \(0<\varepsilon<1\) 时，方程 \(y=x-\varepsilon \sin x\) 可以唯一确定一个函数 \(x=f(y)\) ，它是定义在\((-\infty,+\infty)\) 上的严格单调递增函数．

\subsection{题目4}
4．证明：当 \(a>0\) 时，含参量反常积分 \(\int_{0}^{+\infty} x e^{-x y} \mathrm{~d} y\) 在 \([a,+\infty)\) 上一致收敛，但在 \((0,+\infty)\) 上不一致收敛。

\subsection{题目5}
5．设函数 \(f(x)\) 在（0 ，1］内连续且可导，并且 \(\lim _{x \rightarrow 0^{+}} \sqrt{x} f^{\prime}(x)=a\) ．证明：\(f(x)\) 在（0 ，1］内一致连续．

\subsection{题目6}
6．设 \(n\) 阶实方阵 \(A\) 满足 \(A^{3}=A\) 。证明：若对于任意实向量 \(\alpha\) ，均有 \(\alpha^{T} A^{T} A \alpha \leq \alpha^{T} \alpha\) ，则 \(A\) 是实对称矩阵。

\subsection{题目7}
7．设 \(A, B\) 为 \(n\) 阶矩阵。证明：

\begin{itemize}
  \item[（1）] 若 \(A B=O\) ，则 \(r(A)+r(B) \leq n\) 。
  \item[（2）] 若 \(A^{2}=E\) ，则 \(r(A+E)+r(A-E)=n\) 。
\end{itemize}

\subsection{题目8}
8．设 \(A\) 为复方阵，\(p(x)\) 为 \(E-\bar{A} A\) 的特征多项式，其中 \(\bar{A}\) 是 \(A\) 的共轭矩阵，证明：\(p(x)\) 必为实系数多项式。

\subsection{题目9}
1．一质点在变力 \(\mathbf{F}=\left(x y+e^{-x} \cos y, x^{2}+e^{-x} \sin y\right)\) 的作用下，从原点沿抛物线 \(y=x^{2}\) 运动到点\(A(1,1)\) ，再沿直线运动到点 \(B(0,2)\) ，求该过程中变力所做的功．

\subsection{题目10}
2．当 \(a, b\) 满足何条件时，方程组有解？有解的情况下，求出全部解。

\[
\left\{\begin{array}{l}
a x_{1}+x_{2}+x_{3}=4 \\
x_{1}+b x_{2}+x_{3}=3 \\
x_{1}+2 b x_{2}+x_{3}=4
\end{array}\right.
\]
\end{document}
